\documentclass[10pt,letterpaper,twocolumn]{article}

\usepackage[utf8]{inputenc}
\usepackage[T1]{fontenc}
\usepackage{newtxtext,newtxmath}
\usepackage[scaled=0.92]{helvet}
\usepackage[top=0.76in,bottom=0.82in,left=0.72in,right=0.72in,columnsep=0.23in]{geometry}
\usepackage{microtype}
\usepackage{amsmath,mathtools}
\usepackage{graphicx}
\usepackage{booktabs}
\usepackage{longtable}
\usepackage{array}
\usepackage{enumitem}
\usepackage{titlesec}
\usepackage{xcolor}
\usepackage[hidelinks]{hyperref}
\usepackage{fancyhdr}
\usepackage{setspace}

\setstretch{0.995}
\setlength{\parindent}{0.9em}
\setlength{\parskip}{0pt}
\setlength{\columnsep}{0.23in}
\setlist[itemize]{leftmargin=1.2em,itemsep=0.2em,topsep=0.35em}
\setlist[enumerate]{leftmargin=1.3em,itemsep=0.2em,topsep=0.35em}
\setlength{\abovedisplayskip}{5pt plus 2pt minus 2pt}
\setlength{\belowdisplayskip}{5pt plus 2pt minus 2pt}
\setlength{\abovedisplayshortskip}{4pt plus 2pt minus 2pt}
\setlength{\belowdisplayshortskip}{4pt plus 2pt minus 2pt}
\setlength{\textfloatsep}{8pt plus 2pt minus 2pt}
\setlength{\floatsep}{7pt plus 2pt minus 2pt}
\setlength{\intextsep}{7pt plus 2pt minus 2pt}
\setlength{\emergencystretch}{1.5em}
\raggedbottom

\definecolor{TitleRule}{HTML}{D9DEE7}
\definecolor{TitleNote}{HTML}{4B5563}
\definecolor{HeadingInk}{HTML}{1F2937}
\definecolor{HeadingRule}{HTML}{CBD5E1}

\renewcommand{\thesection}{\arabic{section}}
\renewcommand{\thesubsection}{\thesection.\arabic{subsection}}
\renewcommand{\thesubsubsection}{\thesubsection.\arabic{subsubsection}}
\setcounter{secnumdepth}{2}

\titleformat{\section}
  {\large\sffamily\bfseries\color{HeadingInk}}
  {\thesection}
  {0.65em}
  {}
\titleformat{name=\section,numberless}
  {\large\sffamily\bfseries\color{HeadingInk}}
  {}
  {0pt}
  {}
\titleformat{\subsection}
  {\normalsize\sffamily\bfseries\color{HeadingInk}}
  {\thesubsection}
  {0.55em}
  {}
\titleformat{\subsubsection}
  {\normalsize\sffamily\bfseries\color{HeadingInk}}
  {\thesubsubsection}
  {0.5em}
  {}
\titlespacing*{\section}{0pt}{0.95em}{0.35em}
\titlespacing*{\subsection}{0pt}{0.75em}{0.25em}
\titlespacing*{\subsubsection}{0pt}{0.55em}{0.2em}

\pagestyle{fancy}
\fancyhf{}
\fancyfoot[C]{\small\thepage}
\renewcommand{\headrulewidth}{0pt}

\title{Geodesic Provenance and Structured Source Separation for Photon-Ring Inference}
\author{Jonathan R. Landers}
\date{}

\newcommand{\affiliationnote}{Independent researcher.\\
This note extends the Fourier-domain BHEX analysis in \cite{LandersFourierBHEX} and is inspired by published BHEX work; it is not affiliated with the BHEX collaboration or mission team.}

\makeatletter
\renewcommand{\maketitle}{%
  \twocolumn[{%
    \begin{@twocolumnfalse}
      \vspace*{-1.3em}
      {\LARGE\sffamily\bfseries\raggedright \@title \par}
      \vspace{0.35em}
      {\normalsize\sffamily\raggedright \@author \par}
      \vspace{0.16em}
      {\small\color{TitleNote}\raggedright \affiliationnote \par}
      \vspace{0.65em}
      {\color{TitleRule}\hrule height 0.8pt}
      \vspace{0.8em}
    \end{@twocolumnfalse}
  }]
}
\makeatother

\begin{document}
\maketitle

In a companion note, \emph{A Fourier-Domain Analysis of BHEX for Photon-Ring Inference} \cite{LandersFourierBHEX}, we proposed a structured source-separation framework for interferometric photon-ring inference motivated by the BHEX visibility program \cite{BHEXVision,BHEXDetection}. The central model took the form
\[
y = \alpha g_\theta + q + \varepsilon,
\]
where $g_\theta$ is a photon-ring visibility family, $q$ is structured nuisance, and $\varepsilon$ is noise. The purpose of the present note is to refine that model by adding \emph{geodesic provenance}: observed photons arise only from initial conditions whose null geodesics escape to the observer, while the complementary set is captured by the black hole and never contributes to the data.

The mathematical point is simple. The inverse problem itself does not change, but the nuisance class becomes smaller and more physical. This sharpened provenance-aware formulation leads naturally to improved coherence, mismatch, and stability statements, while remaining close in spirit to the original BHEX-inspired formalism \cite{LandersFourierBHEX,BHEXDetection,BHEXSpin}.

\section{Original formalism and provenance extension}

The original visibility-space model is
\[
y = \alpha g_\theta + q + \varepsilon.
\]
Here $g_\theta$ is the structured photon-ring component, $q$ is nuisance structure such as broader plasma emission, and $\varepsilon$ is noise. In the companion note, this supported both a heuristic ring-oriented estimator and a structured estimator obtained by jointly fitting the ring and nuisance components \cite{LandersFourierBHEX}.

We now add geodesic provenance. Let
\[
z=(p,\xi)
\]
denote a photon initial condition, where $p$ is an emission event and $\xi$ is a null direction. Let $\gamma_z$ be the resulting null geodesic. Define the escape and capture sets
\[
E = \{z:\gamma_z \text{ escapes}\},
\qquad
C = \{z:\gamma_z \text{ is captured}\}.
\]
Only $E$ contributes to observed data. If $f(z)$ is an emission distribution over initial conditions, then the observed image may be written as
\[
I = T_E f,
\]
where $T_E$ is the pushforward map from escaping geodesics to image-plane intensity. The measured visibility is therefore
\[
y = \mathcal F(T_E f) + \varepsilon,
\]
with $\mathcal F$ the image-to-visibility map.

The refinement comes from splitting the escape set as
\[
E = E_{\mathrm{ring}} \,\dot\cup\, E_{\mathrm{bg}},
\]
where $E_{\mathrm{ring}}$ denotes near-critical escaping geodesics that generate the photon ring and $E_{\mathrm{bg}}$ denotes the remaining escaping geodesics. Then
\[
T_E f = T_{E_{\mathrm{ring}}}f + T_{E_{\mathrm{bg}}}f.
\]
Writing
\[
h_\theta := T_{E_{\mathrm{ring}}}f_\theta,
\qquad
b := T_{E_{\mathrm{bg}}}f,
\]
we obtain in image space
\[
I = \alpha h_\theta + b,
\]
and after Fourier transform
\[
g_\theta := \mathcal F h_\theta,
\qquad
q := \mathcal F b,
\]
so that the visibility model again becomes
\[
y = \alpha g_\theta + q + \varepsilon.
\]
Thus the algebraic form is unchanged, but now both signal and nuisance are tagged by geodesic provenance.

\section{Provenance-constrained nuisance classes}

In the companion note, nuisance was modeled abstractly through a subspace or comparison class such as $\mathcal P$ or $Q$ \cite{LandersFourierBHEX}. Here we refine that class by requiring nuisance to arise from non-ring escaping geodesics only. Given an admissible emission family $\mathcal A$, define
\[
Q_{\mathrm{prov}}
:=
\bigl\{\mathcal F T_{E_{\mathrm{bg}}}f : f\in\mathcal A\bigr\}.
\]
Likewise, if $\mathcal P$ denotes the original nuisance subspace in the simpler projection formulation, we write $\mathcal P_{\mathrm{prov}}\subseteq\mathcal P$ for its provenance-restricted analogue.

The physical content is that captured photons never appear in the data, and nuisance can no longer be arbitrary; it must be generated by ordinary escaping geodesics. This is precisely the kind of sharpening suggested by the BHEX perspective, where the photon-ring contribution comes from a distinguished family of strongly lensed trajectories and is separated from broader astrophysical emission by its geometric origin \cite{BHEXDetection,BHEXSpin}.

\section{Coherence improvement under provenance restriction}

In the companion note, one key quantity is the ring--nuisance coherence
\[
\mu(g,\mathcal P)=\frac{\|\Pi_{\mathcal P}g\|}{\|g\|},
\]
where $\Pi_{\mathcal P}$ denotes orthogonal projection onto the nuisance subspace \cite{LandersFourierBHEX}. Provenance restriction improves this quantity immediately.

\par\smallskip\noindent\textbf{Theorem (Monotonicity of coherence under provenance restriction).} Let $\mathcal P_{\mathrm{prov}}\subseteq\mathcal P$ be closed subspaces of the ambient Hilbert space $\mathcal H$, and let $g\in\mathcal H$ be nonzero. Then
\[
\mu(g,\mathcal P_{\mathrm{prov}})\le \mu(g,\mathcal P).
\]
Consequently, restricting nuisance to a provenance-constrained class cannot increase worst-case ring--nuisance overlap.

\par\smallskip\noindent\textbf{Proof.} Since $\mathcal P_{\mathrm{prov}}\subseteq\mathcal P$, orthogonal projection onto $\mathcal P$ gives the best approximation to $g$ from the larger set, while projection onto $\mathcal P_{\mathrm{prov}}$ gives the best approximation from the smaller set. Hence
\[
\|\Pi_{\mathcal P_{\mathrm{prov}}}g\|\le \|\Pi_{\mathcal P}g\|.
\]
Dividing by $\|g\|$ yields the result.

\par\smallskip\noindent\textbf{Corollary.} Any mismatch or stability bound from the companion note controlled by $\mu(g,\mathcal P)$ improves, or at worst remains unchanged, after replacing $\mathcal P$ by $\mathcal P_{\mathrm{prov}}$.

\par\smallskip\noindent\textbf{Equivalent $Q$-formulation.} In the full source-separation model, if ring--nuisance coherence is measured over an admissible nuisance family $Q$ by
\[
\mu(g,Q):=\sup_{q\in Q\setminus\{0\}} \frac{|\langle g,q\rangle|}{\|g\|\,\|q\|},
\]
and if provenance restriction yields a smaller admissible family $Q_{\mathrm{prov}}\subseteq Q$, then
\[
\mu(g,Q_{\mathrm{prov}})\le \mu(g,Q).
\]
The proof is immediate, since the supremum is taken over a smaller set. Thus the same monotonicity holds whether coherence is expressed through projection onto a nuisance subspace or directly over an admissible nuisance class.

\section{Refined mismatch statement}

The companion note proved a mismatch bound between a direct ring-oriented heuristic and a nuisance-aware estimator, controlled by nuisance size, ring--nuisance coherence, and noise \cite{LandersFourierBHEX}. Replacing the original nuisance class by its provenance-restricted version yields the immediate refinement below.

\par\smallskip\noindent\textbf{Lemma (Provenance-refined mismatch).} In the projection-based setting of \cite{LandersFourierBHEX}, replacing $\mathcal P$ by $\mathcal P_{\mathrm{prov}}$ gives
\[
\bigl|\hat\alpha_{\mathrm{heur}}-\hat\alpha_{\mathrm{model}}\bigr|
\le
\mu(g,\mathcal P_{\mathrm{prov}})\frac{\|p\|}{\|g\|}
+
\frac{\|\varepsilon\|}{\|g\|}
+
\frac{\|(I-\Pi_{\mathcal P_{\mathrm{prov}}})\varepsilon\|}{\|(I-\Pi_{\mathcal P_{\mathrm{prov}}})g\|}.
\]
Because $\mu(g,\mathcal P_{\mathrm{prov}})\le \mu(g,\mathcal P)$, this is sharper than the original bound.

The interpretation is simple: the heuristic is not distorted by arbitrary nuisance, only by nuisance that can arise from ordinary escaping geodesics and still overlap the ring template.

\section{Near-critical refinement}

Not all provenance-constrained nuisance is equally harmful. Split the background escape set as
\[
E_{\mathrm{bg}} = E_{\mathrm{near}}\,\dot\cup\,E_{\mathrm{far}},
\]
where $E_{\mathrm{near}}$ consists of escaping background geodesics whose Fourier signatures overlap the oscillatory ring regime, and $E_{\mathrm{far}}$ consists of the remainder. Then
\[
q=q_{\mathrm{near}}+q_{\mathrm{far}}.
\]
In this formulation, the dominant mismatch term is governed more precisely by
\[
\mu(g,Q_{\mathrm{near}})\frac{\|q_{\mathrm{near}}\|}{\|g\|},
\]
rather than by the full nuisance norm. Thus the direct heuristic is destabilized specifically by near-critical contamination, not by all escaping background emission.

\section{Refined identifiability and stability}

The companion note also argued that heuristic readability can fail before true recoverability is lost, and expressed structured recovery through a local identifiability margin \cite{LandersFourierBHEX}. Provenance restriction sharpens that margin by shrinking the admissible nuisance family. Define the provenance-aware identifiability constant by
\[
c_{\mathrm{prov}}
:=
\inf_{q,q'\in Q_{\mathrm{prov}}}
\frac{\|\alpha g_\theta + q - \alpha' g_{\theta'} - q'\|}
{|\alpha-\alpha'| + d(\theta,\theta') + \|q-q'\|}.
\]
Since $Q_{\mathrm{prov}}\subseteq Q$, one has
\[
c_{\mathrm{prov}}\ge c,
\]
where $c$ is the original identifiability margin.

\par\smallskip\noindent\textbf{Proposition (Provenance-improved stability).} Under the provenance-constrained model,
\[
|\hat\alpha-\alpha| + d(\hat\theta,\theta) + \|\hat q-q\|
\le
\frac{C}{c_{\mathrm{prov}}}\eta,
\qquad c_{\mathrm{prov}}\ge c.
\]
Thus restricting nuisance to physically admissible escaping trajectories strengthens the recoverability guarantee.

\section{Discussion}

This extension leaves the inverse problem in the same formal shape as before, but gives it a sharper physical interpretation:
\[
\text{observed signal} = \text{near-critical escaping flow} + 
\]
\[
\text{ordinary escaping flow} + \text{noise}.
\]
From the BHEX viewpoint, this is natural. The long-baseline ring observable is associated with a special class of strongly lensed trajectories, while the broader plasma contribution comes from other escaping emission \cite{BHEXVision,BHEXDetection,BHEXSpin}. The present note simply folds that physical distinction into the abstract source-separation formalism of \cite{LandersFourierBHEX}.

The payoff is mathematical as well as conceptual. Provenance shrinks the nuisance family, lowers coherence, sharpens mismatch bounds, and improves identifiability. In that sense, it provides a clean bridge between the abstract Fourier-domain analysis of the companion note and the geodesic picture underlying BHEX-style photon-ring inference.

\begin{thebibliography}{9}

\bibitem{LandersFourierBHEX}
Jonathan R. Landers,
\emph{A Fourier-Domain Analysis of BHEX for Photon-Ring Inference}, 2026.

\bibitem{BHEXVision}
Black Hole Explorer Collaboration,
\emph{Black Hole Explorer: Motivation and Vision}, 2024.

\bibitem{BHEXDetection}
Black Hole Explorer Collaboration,
\emph{Photon Ring Science, Detection, and Shape Measurement}, 2024.

\bibitem{BHEXSpin}
Black Hole Explorer Collaboration,
\emph{Interferometric Inference of Black Hole Spin from Photon Ring Size and Brightness}, 2025.

\end{thebibliography}

\end{document}
